% Self-contained section fragment; requires amsmath and amssymb.
\providecommand{\file}[1]{\texttt{#1}}
\subsection{A concrete obstruction to the three-residue shortcut}
\label{sec:factor-three-obstruction}

The tempting change from coefficient four to coefficient three is not
merely unsupported by the old decoding proof.  With the unchanged fifth
head it makes the numerical calculus inconsistent.  Thus a new interpretation
of those same numerical axioms cannot repair this particular shortcut.

Write
\[
 p(a,b)=\frac{(a+b)(a+b+1)}2+b,
 \quad I(a,b)=3p(a,b),\quad E(a,b)=3p(a,b)+3,
\]
\[
 U(d,A)=3p(A,d)+2,\qquad N(A)=I(A,1).
\]
Retain modus ponens, generalization from $A$ to $U(d,A)$, the restricted
implicational head
\[
 R(a,b,c,u)=I\bigl(I(I(c,u),a),I(I(a,c),I(b,c))\bigr),
 \qquad u\le a,
\]
and the fifth head
\[
 H_5(d,h)=N\bigl(U(d,N(E(h,d)))\bigr).
\]
No equality-reflexivity axiom, equality-substitution rule, separation gate,
or semantic soundness assumption will be used.

\paragraph{Proposition.}
The resulting numerical calculus derives its designated contradiction,
code $1$, by a finite derivation.  Consequently any faithful exhaustive
proof search for this unchanged calculus halts.

\paragraph{Proof.}
First, $R(0,0,0,0)=0$, so code $0$ is a theorem.  The numerical collision is
\[
                 E(1,0)=6=I(0,1).
\]
The generic restricted-implication derivation gives the usual derived
$K$ and $S$ schemas
\[
 K(a,b)=I(a,I(b,a)),\qquad
 S(a,b,c)=I\bigl(I(a,I(b,c)),I(I(a,b),I(a,c))\bigr).
\]
These derivations still apply at coefficient three.  Specifically,
\file{M6.Logic.RestrictedImplication.k} and
\file{M6.Logic.RestrictedImplication.s} are generic over the implication
operation and comparison relation.  They require only the displayed
restricted head, modus ponens, reflexivity of comparison, and propagation
of a bound into either argument of implication.  On natural numbers with
ordinary $\le$, these bounds follow from
\[
                  a,b\le p(a,b)\le3p(a,b)=I(a,b).
\]
Neither injectivity nor disjoint constructor tags is a hypothesis of those
generic derivations.  The supplied generic Lean source is preserved with
its original attribution in the accompanying certificate directory; no
new Lean compilation is asserted here.

Put $a=6=I(0,1)$.  From $K$ and $S$ derive $I(a,a)$ in the standard way:
apply $S(a,I(a,a),a)$ to $K(a,I(a,a))$ and then to $K(a,a)$.
From $K(0,a)$ and the theorem $0$, derive $I(a,0)$.
Now apply $S(a,0,1)$ first to $I(a,a)=I(a,I(0,1))$ and then to
$I(a,0)$, obtaining
\[
                   I(a,1)=N(6)=87.
\]
Generalization gives
\[
                     U(0,N(6))=11486.
\]
But $H_5(0,1)$ is precisely its negation:
\[
 H_5(0,1)=N(U(0,N(E(1,0))))
        =I(11486,1)=197943987.
\]
One final modus ponens derives $1$.
\hfill$\square$

\paragraph{Executable certificates and scope.}
There are two complementary certificates.  First, \texttt{check.py}
constructs and checks a fourteen-line numeric certificate in
\texttt{proof.json}.  It checks the displayed small code identities and all
instantiated $K/S$ formulas using exact integers.

Second, \texttt{expand\_basis.py} instantiates the retained generic
restricted-implication derivations and expands every use of $K$ and $S$
into primitive rules.  The resulting \texttt{primitive\_proof.json} contains
73,393 proof steps: 34,888 restricted-head instances, 38,503 modus-ponens
steps, one generalization, and one fifth-head instance.  Its conclusions
are represented by a shared directed acyclic graph of 116,171 exact
arithmetic expressions, with maximum expression depth 36.  Sharing avoids
expanding very large numerals and repeated arithmetic subexpressions;
it does not change their values or the calculus.

The separate \texttt{verify\_primitive.py} checker reads only that
certificate.  It neither parses nor trusts the generic Lean source, admits
no $K/S$ rule, and uses no propositional-tautology oracle.  For each
restricted-head side condition it verifies an explicit chain of implication
components, each justified by $a,b\le I(a,b)$.  It checks every
modus-ponens identity by arithmetic-expression congruence, using only the
normalization $I(0,0)=0$; it checks the one generalization and the precise
numeric fifth-head instance.  The final conclusion is literal code $1$.
The final primitive certificate SHA-256 is
\begin{center}\scriptsize\ttfamily
baf316dce45568036fd803382b30ae2669eb504222a39a6be1f3abcad77dcb83.
\end{center}
This supplies a finite primitive-rule refutation with a small independent
checker.  It is not presented as a newly compiled proof-assistant theorem
or as a literal-machine halting trace.

This rules out the unchanged coefficient-three emitter, including any
state saving obtained solely by implementing its three-increment
multiplier more compactly.  A viable three-residue design would need to
change the actual logical interface or emitted heads, and establish a new
calibration.  It cannot retain $R$, $H_5$, generalization and modus ponens
exactly as displayed while preserving consistency.
